% sec-08: the dated record.  Figure 19 (mermaid).
% FINAL stage improvement: Figure 19 gains an eleven-day span marker across the
% contact column, because the compression of the record into eleven days is
% part of what the figure is reporting.
\section{San Diego and the August 2026 Record}

\subsection{Four Contacts, Dated}

Every other section of this document describes a plan. This one describes a
record, and the register changes accordingly: what follows is four contacts and
five responses in eleven days, stated at their actual size.

On \textbf{29 July 2026} the author met one pancreatic cancer surgeon at UC San
Diego Moores Cancer Center. One meeting, one surgeon.

Between \textbf{4 and 8 August 2026} nine of the ten funding applications were
emailed to their recipients, with their attachments and cover text as deposited
\cite{tenapps2026}.

Between \textbf{5 and 7 August 2026} five email responses were received from
those applications.

On \textbf{7 August 2026} an NIH email indicated willingness to explore
additional pancreatic cancer inquiries concerning daraxonrasib.

On \textbf{8 August 2026} UC San Diego Moores Cancer Center emailed an offer of
a two-day seminar together with associated seminar events.

Figure 19 chains the four contacts to what each one actually unlocks, and every
arrow on it ends at the same milestone.

\begin{appfloat}[!tb]
\begin{appfig}[mermaid-type / flowchart]
% Four columns left to right at a 3.85cm pitch, five rows at a 1.42cm pitch.
% The five edges converging on M1 use a symmetric 22/12/0/-12/-22 fan, so the
% five arrivals are at least 3mm apart on M1's west face.
\foreach \i/\y/\sty/\lab in {%
  1/0/mmmid/{2026-07-29\\Moores pancreatic cancer surgeon, one meeting},
  2/-1.42/mmsoft/{2026-08-04 to 08-08\\nine applications emailed},
  3/-2.84/mmsoft/{2026-08-05 to 08-07\\five email responses},
  4/-4.26/mmmid/{2026-08-07\\NIH, explore daraxonrasib inquiries},
  5/-5.68/mmmid/{2026-08-08\\Moores, two-day seminar plus events}}{
  \node[\sty,text width=29mm] (e\i) at (0,\y) {\lab};}
\foreach \i/\y/\lab in {1/0/{A named clinical champion candidate},
  2/-1.42/{Nine review clocks started},
  3/-2.84/{Five contacts and named mechanisms},
  4/-4.26/{A channel for supply and cross reference},
  5/-5.68/{Institutional presence, route to IIT intake}}{
  \node[mmin,text width=27mm] (u\i) at (3.85,\y) {\lab};}
\foreach \i/\y/\lab in {1/0/{Site agreement, IIT slot, budget},
  2/-1.42/{Any award, any obligation},
  3/-2.84/{A program officer commitment},
  4/-4.26/{Letter of authorization, supply agreement},
  5/-5.68/{IRB approval, CTA, theatre time}}{
  \node[mmgray,text width=27mm] (n\i) at (7.70,\y) {\lab};}
\node[mmgoal,text width=30mm,minimum height=16mm] (m1) at (11.90,-2.84)
  {M1 site feasibility executed\\\$24,000, months 1 to 2};
\node[mmlanetitle,anchor=south] at (0,0.92) {Dated contact};
\draw[protoblue,line width=0.7pt] (-1.65,0) -- (-1.65,-5.68);
\draw[protoblue,line width=0.7pt] (-1.75,0) -- (-1.55,0);
\draw[protoblue,line width=0.7pt] (-1.75,-5.68) -- (-1.55,-5.68);
\node[font=\tiny\sffamily\bfseries,text=protoblue,rotate=90,anchor=south] at (-1.80,-2.84)
  {eleven days};
\node[mmlanetitle,anchor=south] at (3.85,0.92) {Unlocked};
\node[mmlanetitle,anchor=south] at (7.70,0.92) {Still not unlocked};
\node[mmlanetitle,anchor=south] at (11.90,0.92) {Terminates at};
\foreach \i in {1,2,3,4,5}{
  \draw[mmedgeb] (e\i) -- (u\i);
  \draw[mmedge] (u\i) -- (n\i);}
\draw[mmedged] (n1.east) to[bend left=22] (m1.north west);
\draw[mmedged] (n2.east) to[bend left=12] (m1.west);
\draw[mmedged] (n3.east) -- (m1.west);
\draw[mmedged] (n4.east) to[bend right=12] (m1.west);
\draw[mmedged] (n5.east) to[bend right=22] (m1.south west);
\node[font=\tiny\itshape,text=protogray,anchor=north west,text width=132mm]
  at (-1.80,-6.85) {Rows 2 and 3 are drawn paler than the other three because they
  record correspondence rather than a meeting, and the distinction has to
  survive without reading the labels. All five arrows terminate at the same
  milestone and not one of the five closes it.};
\end{appfig}
\vspace{-0.65cm}
\figcaption{Figure 19. Four contacts between July 29 and August 8, 2026, and what\\
each one actually unlocks. Every arrow ends at the same milestone,\\
M1, and not one of the four is on its own sufficient to close that.}
\end{appfloat}

\subsection{What Each One Unlocks, and What It Does Not}

The surgeon meeting unlocks a named clinical champion candidate. It does not
unlock a site agreement, a slot in the investigator-initiated trial intake, or a
budget.

The nine emailed applications start nine review clocks. They do not create an
award or an obligation of any kind on any recipient.

The five responses identify five points of contact and five named mechanisms.
They are not a program officer's commitment, and none of the five is a statement
about this trial's merits.

The NIH daraxonrasib email opens a channel for inquiries about supply and
regulatory cross-reference. It does not produce a letter of authorization or a
drug supply agreement.

The Moores seminar offer opens institutional presence and a route toward the
intake committee. It does not produce IRB approval, a clinical trial agreement,
or theatre time.

\subsection{Why the Institution Is the Constraint, Not the Money}

The finding is carried over from the ten-application analysis rather than
re-derived: the programme's schedule is governed by the institution and not by
any funder, and an award arriving in month four cannot begin work until the site
agreement executes \cite{tenapps2026}. That is guard G4 in Figure 7, the only
guard the company cannot satisfy by working harder, and it is halt point H4 in
Figure 15, the earliest of the five at month six.

It is also why the cheapest row in the entire milestone schedule is the first
one. M1 costs \$24,000, which is 1.5 percent of the award, and everything after
it is downstream.

\subsection{The NIH Daraxonrasib Channel}

The 7 August email establishes a channel for inquiries about the agent. It is
worth stating precisely what does not follow from it, because a reader scanning
for traction will be tempted to read more into it than it contains. There is no
letter of authorization. There is no drug supply agreement. There is no
regulatory cross-reference to any existing IND. There has been no communication
of any kind with the agent's developer.

Two rows of the absent register in Table 7 would eventually be closed through
this channel: the drug supply agreement, which blocks M6, and the letter of
authorization, which blocks M5. Revolution Medicines is the only signatory for
both, and no approach has been made.

\subsection{The Seminar, and What It Is For}

The 8 August offer of a two-day seminar with associated events is the most
useful of the four contacts, for a reason that has nothing to do with the
science. It is the route to institutional presence and, through it, to the
investigator-initiated trial intake committee, which is the body that controls
M1.

All four contacts advance M1. None closes it. M1 costs \$24,000 and is the
cheapest row in Table 14, and until it closes nothing else in this document can
begin.
